\chapter{Implementation and Results}
This chapter presents the technical implementation of the AlgoBugs evaluation pipeline and the experimental outcomes from running it across five large language models under three prompting strategies.

\section{Environment Setup}

The AlgoBugs pipeline runs on a personal Linux workstation (Ubuntu 22.04 LTS, Intel Core i5 CPU, 16 GB RAM). All LLM inference is handled remotely through the OpenRouter API; the local machine is responsible only for the Python evaluation engine and g++ compilation. The full evaluation across all 15 model-strategy combinations took approximately 72 hours of wall-clock time, limited by API rate limits rather than local compute.

\subsection{Software Stack}

The evaluation engine (\texttt{evaluate\_models.py}) is written in Python 3.10. The key dependencies are:

\begin{itemize}
    \item \textbf{asyncio / aiohttp}: Drives concurrent API calls. A semaphore caps parallelism at five requests; blocking g++ compilations run in a thread pool via \texttt{asyncio.to\_thread()} so they do not stall the event loop.
    \item \textbf{g++ 11 (C++20)}: Compiles both the buggy and correct solution for every pair using \texttt{-std=c++20 -O2}, matching the Codeforces judge environment.
    \item \textbf{OpenRouter API}: Provides a single endpoint and billing account for all five providers, removing the need for separate API keys per model.
    \item \textbf{Python standard library}: \texttt{json}, \texttt{subprocess}, \texttt{zipfile}, \texttt{pathlib}, and \texttt{re} handle dataset access, compilation, execution, and result serialisation with no additional package dependencies.
\end{itemize}

\subsection{Model Configuration}

Table 4.1 lists the five models evaluated, their OpenRouter identifiers, capability tier, and the number of pairs completed. GPT-5.1 stopped at 870 pairs due to a project credit limit; all FER values for that model are computed over its 870-pair sample only.

\begin{table}[h!]
    \centering
    \caption{Models and Configuration}
    \label{tab:models}
    \begin{tabular}{lllc}
        \toprule
        \textbf{Model} & \textbf{OpenRouter ID} & \textbf{Tier} & \textbf{Pairs} \\
        \midrule
        Gemini 2.0 Flash   & \texttt{google/gemini-2.0-flash-001}           & Mid-tier        & 1,000 \\
        LLaMA 3.3 70B      & \texttt{meta-llama/llama-3.3-70b-instruct}     & Open-source     & 1,000 \\
        Qwen3 Coder 30B    & \texttt{qwen/qwen3-coder-30b-a3b-instruct}     & Code-specialist & 1,000 \\
        DeepSeek V3        & \texttt{deepseek/deepseek-chat}                & Open-source     & 1,000 \\
        GPT-5.1            & \texttt{openai/gpt-5.1}                        & Frontier        &   870 \\
        \bottomrule
    \end{tabular}
\end{table}

\subsection{Pipeline Architecture}

Figure 4.1 shows the runtime flow for a single pair. Each pair passes through prompt construction, LLM inference, test extraction, compilation of both solutions, differential execution, and verdict assignment.

\begin{figure}[h!]
    \centering
    \begin{tikzpicture}[
        box/.style={rectangle, rounded corners=3pt, draw=black!70, fill=blue!10,
                    text width=2.5cm, align=center, font=\small,
                    minimum height=0.9cm, inner sep=4pt},
        vbox/.style={rectangle, rounded corners=3pt, draw=black!70, fill=green!15,
                     text width=2.5cm, align=center, font=\small,
                     minimum height=0.9cm, inner sep=4pt},
        arr/.style={->, >=Stealth, draw=black!60, thick}
    ]
    \node[box]  (A) at (0.0,  0)    {Dataset Pair\\(problem + code)};
    \node[box]  (B) at (3.2,  0)    {Prompt\\Construction};
    \node[box]  (C) at (6.4,  0)    {LLM via\\OpenRouter};
    \node[box]  (D) at (9.6,  0)    {Test Input\\Extracted};
    \node[box]  (E) at (9.6, -2.3)  {g++ Compile\\Both Solutions};
    \node[box]  (F) at (6.4, -2.3)  {Execute on\\Test Input};
    \node[box]  (G) at (3.2, -2.3)  {Compare\\Outputs};
    \node[vbox] (H) at (0.0, -2.3)  {Verdict\\(FER)};

    \draw[arr] (A) -- (B);
    \draw[arr] (B) -- (C);
    \draw[arr] (C) -- (D);
    \draw[arr] (D) -- (E);
    \draw[arr] (E) -- (F);
    \draw[arr] (F) -- (G);
    \draw[arr] (G) -- (H);

    \node[draw=black!40, dashed, rounded corners=2pt, fill=yellow!8,
          text width=2.2cm, align=center, font=\scriptsize, inner sep=3pt]
          (strat) at (4.8, 0.95) {ZS / CoT / FS};
    \draw[->, dashed, draw=black!40] (strat.south) -- ++(0,-0.35);
    \end{tikzpicture}
    \caption{AlgoBugs Evaluation Pipeline Runtime Flow}
    \label{fig:pipeline}
\end{figure}

\subsection{Verdict Classification and FER Denominator}

Each evaluated pair receives exactly one of seven verdicts:

\begin{itemize}
    \item \textbf{BUG\_EXPOSED}: the buggy solution and the correct solution produce different output on the generated test input. This is the success case.
    \item \textbf{NOT\_EXPOSED}: both solutions produce identical output. The test is valid but does not reveal the bug.
    \item \textbf{COMPILE\_ERROR\_BUGGY} / \textbf{COMPILE\_ERROR\_CORRECT}: one solution fails g++ compilation, making the pair unscoreable.
    \item \textbf{NO\_TEST\_GENERATED}: the model response contains no extractable test input.
    \item \textbf{INVALID\_TEST}: a test input is extracted but causes a runtime error on both solutions.
    \item \textbf{ERROR}: an API or execution failure not attributable to model output.
\end{itemize}

The Fault Exposure Rate uses only BUG\_EXPOSED and NOT\_EXPOSED in its denominator:

\begin{equation}
    \text{FER} = \frac{|\text{BUG\_EXPOSED}|}{|\text{BUG\_EXPOSED}| + |\text{NOT\_EXPOSED}|} \times 100\%
\end{equation}

Compile errors and format failures are excluded from both numerator and denominator. FER therefore measures fault-exposure performance on the subset of pairs where the model did produce an executable test, without penalising or inflating the score for format failures.

\subsection{Rate Limiting and Fault Tolerance}

The engine applies exponential backoff (base 2 seconds) on HTTP 429 and server-error responses. A circuit breaker pauses all requests to the affected model for 60 seconds after five consecutive failures. Results are checkpointed to disk every ten pairs, so a network interruption or credit-limit stop does not discard completed work. GPT-5.1 was stopped by the credit limit mid-run; the 870 completed pairs were fully preserved in the checkpoint file.

\section{Testing and Evaluation}

\subsection{Overall FER by Model and Strategy}

Table 4.2 shows the global FER for each model under all three prompting strategies. The range across all 15 model-strategy combinations is 10.8\% to 21.2\%.

\begin{table}[h!]
    \centering
    \caption{Overall Fault Exposure Rate (\%) by Model and Prompting Strategy}
    \label{tab:overall_fer}
    \begin{tabular}{lcccc}
        \toprule
        \textbf{Model} & \textbf{Zero-Shot} & \textbf{Chain-of-Thought} & \textbf{Few-Shot} & \textbf{n} \\
        \midrule
        Gemini 2.0 Flash   & 15.6 & 15.4 & 14.9 & 1,000 \\
        LLaMA 3.3 70B      & 17.8 & 20.0 & 21.2 & 1,000 \\
        Qwen3 Coder 30B    & 20.2 & 17.4 & 19.3 & 1,000 \\
        DeepSeek V3        & 14.2 & 10.8 & 14.5 & 1,000 \\
        GPT-5.1            & 14.6 & 15.1 & 20.2 &   870 \\
        \bottomrule
    \end{tabular}
\end{table}

Figure 4.2 visualises the same data as a grouped bar chart, making cross-strategy differences visible per model.

\begin{figure}[h!]
    \centering
    \begin{tikzpicture}
    \begin{axis}[
        ybar,
        bar width=0.22cm,
        width=\linewidth,
        height=7.5cm,
        enlarge x limits=0.12,
        ylabel={FER (\%)},
        symbolic x coords={Gemini 2.0 Flash, LLaMA 3.3 70B, Qwen3 Coder 30B, DeepSeek V3, GPT-5.1},
        xtick=data,
        x tick label style={font=\small, rotate=18, anchor=east},
        ymin=0, ymax=27,
        legend style={at={(0.5,1.03)}, anchor=south, legend columns=3, font=\small},
        ymajorgrids=true,
        grid style={dashed, draw=black!20},
    ]
    \addplot[fill=blue!60, draw=blue!80] coordinates {
        (Gemini 2.0 Flash, 15.6)
        (LLaMA 3.3 70B,    17.8)
        (Qwen3 Coder 30B,  20.2)
        (DeepSeek V3,      14.2)
        (GPT-5.1,          14.6)
    };
    \addplot[fill=red!55, draw=red!75] coordinates {
        (Gemini 2.0 Flash, 15.4)
        (LLaMA 3.3 70B,    20.0)
        (Qwen3 Coder 30B,  17.4)
        (DeepSeek V3,      10.8)
        (GPT-5.1,          15.1)
    };
    \addplot[fill=green!55, draw=green!75] coordinates {
        (Gemini 2.0 Flash, 14.9)
        (LLaMA 3.3 70B,    21.2)
        (Qwen3 Coder 30B,  19.3)
        (DeepSeek V3,      14.5)
        (GPT-5.1,          20.2)
    };
    \legend{Zero-Shot, Chain-of-Thought, Few-Shot}
    \end{axis}
    \end{tikzpicture}
    \caption{Fault Exposure Rate by Model and Prompting Strategy}
    \label{fig:overall_fer_chart}
\end{figure}

Qwen3 Coder 30B leads under zero-shot at 20.2\%, with LLaMA 3.3 70B at 17.8\%. GPT-5.1's zero-shot value of 14.6\% is over a smaller sample (870 pairs), so direct comparison with the other four requires some care. DeepSeek V3 shows the widest swing across strategies: 14.2\% zero-shot falls to 10.8\% under chain-of-thought before recovering to 14.5\% with few-shot.

\subsection{Per-Category FER (Zero-Shot)}

Table 4.3 breaks down zero-shot FER by bug category for all five models. This is the diagnostic core of AlgoBugs: it shows where models succeed and where they consistently fail, rather than averaging over categories with very different difficulty.

\begin{table}[h!]
    \centering
    \caption{Per-Category Zero-Shot FER (\%) Across All Five Models}
    \label{tab:per_category_fer}
    \resizebox{\linewidth}{!}{%
    \begin{tabular}{lrrrrr}
        \toprule
        \textbf{Category} & \textbf{Gemini} & \textbf{LLaMA} & \textbf{Qwen} & \textbf{DeepSeek} & \textbf{GPT-5.1} \\
        \midrule
        T1: Integer Overflow     &  1.7 &  6.6 &  8.2 &  1.6 &  3.3 \\
        T2: Modular Arithmetic   &  2.9 & 11.8 & 13.9 &  5.4 &  8.1 \\
        T3: Off-by-One           & 29.5 & 22.8 & 21.8 & 26.3 & 27.7 \\
        T4: Wrong Conditional    & 16.7 & 17.5 & 18.8 & 14.6 & 16.8 \\
        T5: Algorithmic Flaw     & 13.7 & 19.5 & 21.7 & 14.0 & 11.7 \\
        T6: State/Initialization & 10.7 & 10.8 & 21.6 & 11.0 &  8.0 \\
        T7: Corner Case          & 14.6 & 24.4 & 20.5 & 15.2 &  2.2 \\
        T8: I/O Format           & 15.2 & 19.6 & 35.6 &  4.3 & 20.8 \\
        \bottomrule
    \end{tabular}%
    }
\end{table}

Figure 4.3 shows the full per-model, per-category FER as a heat map. Each cell gives the zero-shot FER for one model--category pair; shade intensity is proportional to the value, with darker cells indicating higher fault-exposure rates.

\begin{figure}[h!]
\centering
\setlength{\tabcolsep}{3pt}
\renewcommand{\arraystretch}{1.35}
{\small
\begin{tabular}{l|c|c|c|c|c|c|c|c|}
    \cline{2-9}
    & \rotatebox{55}{\textbf{\scriptsize T1: Int.\ Overflow}\hspace{1mm}}
    & \rotatebox{55}{\textbf{\scriptsize T2: Mod.\ Arith.}\hspace{1mm}}
    & \rotatebox{55}{\textbf{\scriptsize T3: Off-by-One}\hspace{1mm}}
    & \rotatebox{55}{\textbf{\scriptsize T4: Wrong Cond.}\hspace{1mm}}
    & \rotatebox{55}{\textbf{\scriptsize T5: Algo.\ Flaw}\hspace{1mm}}
    & \rotatebox{55}{\textbf{\scriptsize T6: State/Init.}\hspace{1mm}}
    & \rotatebox{55}{\textbf{\scriptsize T7: Corner Case}\hspace{1mm}}
    & \rotatebox{55}{\textbf{\scriptsize T8: I/O Format}\hspace{1mm}} \\
    \hline
    \scriptsize DeepSeek V3
        & \cellcolor{blue!4}\scriptsize 1.6
        & \cellcolor{blue!12}\scriptsize 5.4
        & \cellcolor{blue!59}\textcolor{white}{\scriptsize 26.3}
        & \cellcolor{blue!33}\scriptsize 14.6
        & \cellcolor{blue!31}\scriptsize 14.0
        & \cellcolor{blue!25}\scriptsize 11.0
        & \cellcolor{blue!34}\scriptsize 15.2
        & \cellcolor{blue!10}\scriptsize 4.3 \\\hline
    \scriptsize Gemini Flash
        & \cellcolor{blue!4}\scriptsize 1.7
        & \cellcolor{blue!7}\scriptsize 2.9
        & \cellcolor{blue!66}\textcolor{white}{\scriptsize 29.5}
        & \cellcolor{blue!38}\scriptsize 16.7
        & \cellcolor{blue!31}\scriptsize 13.7
        & \cellcolor{blue!24}\scriptsize 10.7
        & \cellcolor{blue!33}\scriptsize 14.6
        & \cellcolor{blue!34}\scriptsize 15.2 \\\hline
    \scriptsize LLaMA 70B
        & \cellcolor{blue!15}\scriptsize 6.6
        & \cellcolor{blue!27}\scriptsize 11.8
        & \cellcolor{blue!51}\textcolor{white}{\scriptsize 22.8}
        & \cellcolor{blue!39}\scriptsize 17.5
        & \cellcolor{blue!44}\scriptsize 19.5
        & \cellcolor{blue!24}\scriptsize 10.8
        & \cellcolor{blue!55}\textcolor{white}{\scriptsize 24.4}
        & \cellcolor{blue!44}\scriptsize 19.6 \\\hline
    \scriptsize GPT-5.1
        & \cellcolor{blue!7}\scriptsize 3.3
        & \cellcolor{blue!18}\scriptsize 8.1
        & \cellcolor{blue!62}\textcolor{white}{\scriptsize 27.7}
        & \cellcolor{blue!38}\scriptsize 16.8
        & \cellcolor{blue!26}\scriptsize 11.7
        & \cellcolor{blue!18}\scriptsize 8.0
        & \cellcolor{blue!5}\scriptsize 2.2
        & \cellcolor{blue!47}\textcolor{white}{\scriptsize 20.8} \\\hline
    \scriptsize Qwen3 Coder
        & \cellcolor{blue!18}\scriptsize 8.2
        & \cellcolor{blue!31}\scriptsize 13.9
        & \cellcolor{blue!49}\textcolor{white}{\scriptsize 21.8}
        & \cellcolor{blue!42}\scriptsize 18.8
        & \cellcolor{blue!49}\textcolor{white}{\scriptsize 21.7}
        & \cellcolor{blue!49}\textcolor{white}{\scriptsize 21.6}
        & \cellcolor{blue!46}\textcolor{white}{\scriptsize 20.5}
        & \cellcolor{blue!80}\textcolor{white}{\scriptsize 35.6} \\\hline
\end{tabular}}
\caption{Per-category zero-shot FER (\%) for all five models. Cell shade (white-to-blue) is proportional to FER value; white text is used where the background is too dark for black text. T3 is consistently the darkest column; T1 and T2 are consistently the lightest, confirming the category-level pattern across all models.}
\label{fig:heatmap}
\end{figure}

Figure 4.4 shows the average FER per category across all five models, ordered from most to least exposed.

\begin{figure}[h!]
    \centering
    \begin{tikzpicture}
    \begin{axis}[
        xbar,
        bar width=0.42cm,
        width=0.88\linewidth,
        height=8.5cm,
        xlabel={Average FER (\%)},
        symbolic y coords={
            T1: Int Overflow,
            T2: Mod Arith,
            T6: State Init,
            T7: Corner Case,
            T5: Algo Flaw,
            T4: Wrong Cond,
            T8: IO Format,
            T3: Off-by-One
        },
        ytick=data,
        y tick label style={font=\small},
        xmin=0, xmax=34,
        nodes near coords,
        nodes near coords align={horizontal},
        every node near coord/.append style={font=\scriptsize, xshift=2pt},
        xmajorgrids=true,
        grid style={dashed, draw=black!20},
        fill=blue!55,
        draw=blue!70,
    ]
    \addplot coordinates {
        (4.3,T1: Int Overflow)
        (8.4,T2: Mod Arith)
        (12.4,T6: State Init)
        (15.4,T7: Corner Case)
        (16.1,T5: Algo Flaw)
        (16.9,T4: Wrong Cond)
        (19.1,T8: IO Format)
        (25.6,T3: Off-by-One)
    };
    \end{axis}
    \end{tikzpicture}
    \caption{Average Per-Category FER (\%) Across All Five Models (Zero-Shot). T3 (Off-by-One) leads at 25.6\%; T1 (Integer Overflow) trails at 4.3\%.}
    \label{fig:per_category_avg}
\end{figure}

T3 (Off-by-One) stands alone at 25.6\% average, 6.5 percentage points above the next category (T8: I/O Format at 19.1\%). T1 (Integer Overflow) at 4.3\% and T2 (Modular Arithmetic) at 8.4\% are the weakest categories by a clear margin.

\subsection{Prompting Strategy Effects}

Figure 4.5 shows the change in FER relative to the zero-shot baseline for each model under chain-of-thought and few-shot prompting. Positive values indicate improvement; negative values indicate decline.

\begin{figure}[h!]
    \centering
    \begin{tikzpicture}
    \begin{axis}[
        ybar,
        bar width=0.28cm,
        width=0.90\linewidth,
        height=6.5cm,
        ylabel={$\Delta$ FER (pp) from zero-shot baseline},
        ylabel style={font=\small},
        symbolic x coords={Gemini Flash, LLaMA 70B, Qwen3 Coder, DeepSeek V3, GPT-5.1},
        xtick=data,
        x tick label style={rotate=18, anchor=east, font=\small},
        ymin=-5.5, ymax=8.0,
        ytick={-4,-2,0,2,4,6},
        enlarge x limits=0.12,
        legend style={at={(0.98,0.98)}, anchor=north east, font=\small},
        ymajorgrids=true,
        grid style={dashed, draw=black!15},
        extra y ticks={0},
        extra y tick labels={},
        extra tick style={grid=major, grid style={black, thick, dashed}},
    ]
    \addplot+[fill=blue!50, draw=blue!65] coordinates {
        (Gemini Flash,-0.2)(LLaMA 70B,2.2)(Qwen3 Coder,-2.8)(DeepSeek V3,-3.4)(GPT-5.1,0.5)
    };
    \addplot+[fill=orange!65, draw=orange!80] coordinates {
        (Gemini Flash,-0.7)(LLaMA 70B,3.4)(Qwen3 Coder,-0.9)(DeepSeek V3,0.3)(GPT-5.1,5.6)
    };
    \legend{Chain-of-Thought $\Delta$, Few-Shot $\Delta$}
    \end{axis}
    \end{tikzpicture}
    \caption{Change in FER (percentage points) from the zero-shot baseline under chain-of-thought (blue) and few-shot (orange) prompting. The dashed line at zero marks the baseline. DeepSeek V3 shows the sharpest CoT decline ($-$3.4~pp); GPT-5.1 shows the largest few-shot gain ($+$5.6~pp).}
    \label{fig:strategy_delta}
\end{figure}

Three patterns emerge from comparing strategies in Table 4.2.

Chain-of-thought prompting produced no consistent improvement. Gemini's FER moved from 15.6\% (zero-shot) to 15.4\% (chain-of-thought), effectively unchanged. DeepSeek V3 dropped by 3.4 percentage points from 14.2\% to 10.8\%, which is the largest strategy-induced change in the entire dataset, and it is a decline. Only LLaMA improved meaningfully under chain-of-thought (+2.2 percentage points).

Few-shot prompting produced model-specific improvements. LLaMA 3.3 70B gained 3.4 percentage points over zero-shot (17.8\% to 21.2\%). GPT-5.1 gained 5.6 percentage points (14.6\% to 20.2\%). The other three models changed by less than 1 percentage point in either direction under few-shot.

One mechanical observation about few-shot: the three demonstration examples anchored the output format precisely. Under zero-shot and chain-of-thought, a small number of responses per model produced no extractable test input. With three worked examples in the prompt, this count fell to zero for all five models. The demonstrations eliminated output format ambiguity even in cases where they did not improve the quality of the test itself.

\section{Results and Discussion}

\subsection{Off-by-One Bugs Are Consistently Exposed}

T3 (Off-by-One) is the most exposed category across all five models at an average of 25.6\% under zero-shot. Every model exceeds 20\%, with Gemini reaching 29.5\% and GPT-5.1 at 27.7\%. Off-by-one bugs trigger on inputs near array boundaries or loop termination conditions. Generating values like $n$, $n-1$, or $n+1$ for an array size or loop bound exercises these cases, and the models appear to infer this pattern directly from reading C++ loop and index expressions. The result is structural: the class of triggering inputs for T3 aligns with the inputs a model naturally generates when it reads a for-loop or array access, so the high FER here reflects pattern recognition, not genuine fault-targeted reasoning.

\subsection{Integer Overflow and Modular Arithmetic Resist All Models}

T1 (Integer Overflow) is the hardest category at 4.3\% average. The best individual result is Qwen3 Coder 30B at 8.2\%; both Gemini and DeepSeek fall below 2\%. T2 (Modular Arithmetic) is nearly as hard at 8.4\% average. Exposing an integer overflow requires the model to identify the specific arithmetic expression that overflows and then choose an input value that satisfies the stated problem constraints while causing the intermediate computation to exceed 32-bit or 64-bit bounds. That requires fault-targeted reasoning about the arithmetic, not code structure pattern-matching. None of the five models did this reliably, regardless of prompting strategy.

\subsection{Chain-of-Thought Provides No Consistent Benefit}

Chain-of-thought prompting was designed to scaffold reasoning: the model identifies the bug, reasons about its trigger condition, and then generates a test. In practice, no consistent improvement resulted. Three of the five models changed by less than one percentage point. DeepSeek V3 declined from 14.2\% to 10.8\%, the sharpest drop in the dataset. One likely cause: the CoT template asked the model to reason about what the bug does before generating a test, and DeepSeek appears to have allocated its context budget to verbose bug descriptions without translating that analysis into more precise numerical test values. Describing what a bug does and generating the specific input that exposes it are different tasks, and CoT as applied here did not bridge that gap.

\subsection{Few-Shot Prompting Is Model-Specific, Not Universally Effective}

The three synthetic few-shot demonstrations (T1: integer overflow, T3: off-by-one, T4: wrong conditional) were drawn from outside the AlgoBugs dataset, so they contain no direct category-level information about most test items. LLaMA 3.3 70B gained 3.4 percentage points and GPT-5.1 gained 5.6 percentage points over their respective zero-shot baselines. Gemini, DeepSeek, and Qwen each changed by less than 1 percentage point.

The asymmetry argues against a simple reading that demonstrations help. A sharper reading is that few-shot demonstrations transferred a reasoning pattern to the two models that were sensitive to task-format guidance, while the remaining three either already applied a similar strategy (Qwen3 Coder 30B at 20.2\% zero-shot was the highest baseline) or processed the demonstrations without adopting the fault-exposure reasoning structure.

\subsection{FER Ceiling Reflects Specification-Valid, Not Fault-Targeted, Generation}

The highest FER in the entire experiment is 21.2\% (LLaMA 3.3 70B, few-shot). Across all 15 model-strategy combinations, the mean FER is approximately 16.5\%. These numbers point to the same gap: all five models generate specification-valid tests at a high rate, but only a fraction of those tests happen to trigger the planted bug.

A specification-valid test satisfies the declared constraints and does not crash either solution. Every model achieves this on the vast majority of pairs. A fault-targeted test additionally aims at the specific computation the bug corrupts. For T3 bugs, the triggering input is near a boundary, which is where models naturally look. For T1 bugs, the triggering input requires a precise large value that overflows an intermediate product, which models consistently miss. That difference explains why T3 FER is three to six times higher than T1 across the board, and why the ceiling of roughly 21\% persists regardless of which model or strategy is used. The 21\% ceiling is not a property of the evaluation setup; it is the empirical answer to the research question at the time of writing.

\section{Summary}

This chapter described the AlgoBugs evaluation pipeline and reported results across 5 models, 3 prompting strategies, and 8 bug categories. The pipeline used asynchronous Python with local g++ execution for differential testing and OpenRouter for LLM access. GPT-5.1 completed 870 of the 1,000 pairs due to a credit limit; all other models completed the full set.

FER ranged from 10.8\% to 21.2\% across all 15 combinations. T3 (Off-by-One) was the most exposed category at a 25.6\% average; T1 (Integer Overflow) at 4.3\% was the hardest. Chain-of-thought prompting showed no consistent improvement and lowered DeepSeek V3's FER by 3.4 percentage points. Few-shot prompting improved LLaMA 3.3 70B (+3.4 pp) and GPT-5.1 (+5.6 pp) but left the other three models unchanged. The consistent ceiling near 21\% reflects specification-valid rather than fault-targeted test generation across all current models tested. Chapter 5 covers the ethical and professional standards that governed the research.
